\documentclass{article}
\usepackage{amsmath, amssymb, amsthm}
\usepackage{hyperref}
\usepackage{geometry}
\geometry{margin=1in}

\title{Erd\H{o}s Problem \#1020}
\date{Last updated: 2025-09-12}
\author{Source: erdosproblems.com}

\begin{document}
\maketitle

\noindent\textbf{Status:} FALSIFIABLE \quad \textbf{No prize}

\noindent\textbf{Tags:} graph theory, hypergraphs

\medskip
\noindent\textbf{Note:} Erdős matching conjecture

\medskip

\noindent\textbf{Problem Statement:}

Let $f(n;r,k)$ be the maximal number of edges in an $r$-uniform hypergraph which contains no set of $k$ many independent edges. For all $r\geq 3$,\[f(n;r,k)=\max\left(\binom{rk-1}{r}, \binom{n}{r}-\binom{n-k+1}{r}\right).\]




Erd\H{o}s and Gallai \cite{ErGa59} proved this is true when $r=2$ (when $r=2$ this also follows from the Erd\H{o s-Ko-Rado theorem}). The conjectured form of $f(n;r,k)$ is the best possible, as witnessed by two examples: all $r$-edges on a set of $rk-1$ many vertices, and all edges on a set of $n$ vertices which contain at least one element of a fixed set of $k-1$ vertices. Frankl \cite{Fr87} proved $f(n;r,k) \leq (k-1)\binom{n-1}{r-1}$. This is sometimes known as the Erd\H{o}s matching conjecture. Note that the second term in the maximum dominates when $n\geq (r+1)k$. There are many partial results towards this, establishing the conjecture in different ranges. These can be separated into two regimes. For small $n$: {UL} {LI}The conjecture is trivially true if $n&lt;kr$.{/LI} {LI}Kleitman \cite{Kl68} when $n=kr$.{/LI} {LI}Frankl \cite{Fr17} when\[kr \leq n\leq k\left(r+\frac{1}{2r^{2r+1}}\right).\]{/LI} {LI}Kolupaev and Kupavskii \cite{KoKu23} when $r\geq 5$, $k&gt;101r^3$, and\[kr \leq n &lt; k\left(r+\frac{1}{100r}\right).\]{/LI} {/UL} For large $n$: {UL} {LI}Erd\H{o}s \cite{Er65d} when $n&gt;kc_r$ (where $c_r$ depends on $r$ in some unspecified fashion).{/LI} {LI} Frankl and F\"{u}redi \cite{Fr87} when $n&gt;100 k^2r$.{/LI} {LI} Bollob\'{a}s, Daykin, and Erd\H{o}s \cite{BDE76} when $n\geq 2kr^3$.{/LI} {LI} Frankl, R\"{o}dl, and Ruci\'{n}ski \cite{FRR12} when $r=3$ and $n\geq 4k$.{/LI} {LI} Huang, Loh, and Sudakov \cite{HLS12} when $n\geq 3kr^2$.{/LI} {LI} Frankl, Luczak, and Mieczkowska \cite{FLM12} when $n&gt; 2k\frac{r^2}{\log r}$.{/LI} {LI} Luczak and Mieczkowska \cite{LuMi14} when $r=3$ (for all $k$).{/LI} {/UL}




References

[BDE76] Bollob\'{a}s, B. and Daykin, D. E. and Erd\H{o}s, P., Sets of independent edges of a hypergraph . Quart. J. Math. Oxford Ser. (2) (1976), 25--32.

[Er65d] Erd\H{o}s, P., A problem on independent {$r$}-tuples . Ann. Univ. Sci. Budapest. E\"otv\"os Sect. Math. (1965), 93--95.

[ErGa59] Erd\H{o}s, P. and Gallai, T., On maximal paths and circuits of graphs . Acta Math. Acad. Sci. Hungar. (1959), 337-356 (unbound insert).

[FLM12] Frankl, Peter and \L uczak, Tomasz and Mieczkowska, Katarzyna, On matchings in hypergraphs . Electron. J. Combin. (2012), Paper 42, 5.

[FRR12] Frankl, Peter and R\"odl, Vojtech and Ruci\'nski, Andrzej, On the maximum number of edges in a triple system not containing a disjoint family of a given size . Combin. Probab. Comput. (2012), 141--148.

[Fr17] Frankl, Peter, Proof of the {E}rd\H{o}s matching conjecture in a new range . Israel J. Math. (2017), 421--430.

[Fr87] Frankl, Peter, The shifting technique in extremal set theory . (1987), 81--110.

[HLS12] Huang, Hao and Loh, Po-Shen and Sudakov, Benny, The size of a hypergraph and its matching number . Combin. Probab. Comput. (2012), 442--450.

[Kl68] Kleitman, Daniel J., Maximal number of subsets of a finite set no {$k$} of which are pairwise disjoint . J. Combinatorial Theory (1968), 157--163.

[KoKu23] Kolupaev, Dmitriy and Kupavskii, Andrey, Erd\H{o}s matching conjecture for almost perfect matchings . Discrete Math. (2023), Paper No. 113304, 9.

[LuMi14] \L uczak, Tomasz and Mieczkowska, Katarzyna, On {E}rd\H{o}s' extremal problem on matchings in hypergraphs . J. Combin. Theory Ser. A (2014), 178--194.

\bigskip
\noindent\small{Source: \url{https://www.erdosproblems.com/1020}}
\end{document}
